\section{Introduction}
\label{sec:introduction}

Congressional redistricting in the United States is not governed by a single national standard. Beyond the federal floor requirements established by \textit{Wesberry v.\ Sanders}\footnote{\textit{Wesberry v.\ Sanders}, 376 U.S. 1 (1964) (requiring congressional districts to be as equal in population as practicable).} and the Voting Rights Act,\footnote{52 U.S.C. \S\,10301 (2018).} each of the 50 states imposes its own redistricting criteria through constitutions, statutes, and court decisions. Some states require nothing more than population equality. Others mandate independent commissions, partisan fairness standards, community-of-interest preservation, county-integrity protections, and explicit compactness requirements. The result is a legal landscape of extraordinary heterogeneity: a congressional map that is fully lawful in Texas would be constitutionally infirm in California, and a map drawn to California's commission standards would fail to satisfy Iowa's contiguity-and-compactness statute.

This heterogeneity poses an immediate challenge for any algorithmic redistricting system aspiring to national applicability. An algorithm that draws maps using a fixed set of parameters will satisfy some states' criteria and fail others'. Yet the administrative appeal of algorithmic redistricting---reproducibility, resistance to partisan manipulation, transparent criteria---is greatest when the algorithm can be applied uniformly across the country, adapting to state-specific legal requirements through parameter configuration rather than through bespoke code. If each state requires a separately designed algorithm, the reproducibility advantages of the algorithmic approach are substantially diminished.

This paper argues that a single algorithmic architecture, parameterized through per-state YAML configuration files, can satisfy the dominant legal criteria across all five categories of state redistricting regimes. The argument has three parts. First, we develop a five-type taxonomy of state redistricting regimes based on the nature and stringency of their legal criteria: \textit{permissive} (minimal criteria beyond federal floor), \textit{criteria-based} (specific mandatory criteria without commission), \textit{commission} (independent commission administering specified criteria), \textit{hybrid} (legislative or commission process with court oversight), and \textit{restrictive} (most stringent criteria including partisan fairness requirements). Second, we show how each category of legal criterion maps onto a concrete algorithmic parameter: compactness requirements onto edge-weight scaling factors, county-preservation requirements onto geographic weight overrides, communities-of-interest requirements onto custom edge-weight matrices, and partisan-neutrality requirements onto the algorithm's structural architecture. Third, we develop model YAML configurations for one representative state in each type, demonstrating the concrete parameterization.

The analysis reveals two striking findings. First, algorithmic redistricting is more legally adaptable than previously recognized: the same graph-partitioning framework can satisfy criteria as different as Texas's minimal requirements and California's commission standards through parameter adjustment alone. Second, the algorithm's structural partisan blindness---its architectural inability to read partisan data---satisfies partisan-neutrality requirements more completely than any human-administered process, because it prevents not merely partisan intent but partisan access. This distinction matters for states like Florida, which constitutionally prohibits drawing maps that ``favor or disfavor a political party,'' and for Arizona, whose commission is statutorily required to maintain partisan fairness. A human mapmaker who promises not to consider partisan data must be trusted on their word; an algorithm that architecturally cannot receive partisan data satisfies the requirement by design.

The paper proceeds as follows. Section~\ref{sec:classification} presents the five-type taxonomy and classifies all 50 states. Section~\ref{sec:algorithm} maps each legal criterion category to its algorithmic parameter. Section~\ref{sec:configs} presents model YAML configurations for each state type. Section~\ref{sec:legal} analyzes which criteria are constitutionally mandated versus merely discretionary. Section~\ref{sec:conclusion} concludes.

\subsection{Scope and Limitations}

This analysis focuses on congressional redistricting criteria, not state legislative criteria, which are separately analyzed in a companion paper.\footnote{See Della-Libera, \textit{Voting Rights Act Compliance for State Legislative Redistricting: The 42\% Threshold at Chamber Scale} (2026) (F.6).} We address 50-state variation in the criteria governing congressional maps: the number of districts, the population equality standard, and any additional criteria imposed by state law. We treat the federal floor---population equality and VRA compliance---as fixed and analyze how state-specific criteria layer on top of it.

A further limitation: our taxonomy is necessarily reductive. Many states' redistricting criteria are the product of decades of litigation, advisory opinions, and administrative practice that a taxonomy cannot fully capture. We classify states by their dominant legal criterion and note important exceptions. Readers working on specific state implementations should consult primary sources and current case law.
